\documentclass[12pt,a4paper]{report}

%% ── Packages ──────────────────────────────────────────────────────────────────
\usepackage[a4paper, margin=2.5cm]{geometry}
\usepackage{fontenc}
\usepackage{inputenc}
\usepackage{microtype}
\usepackage{lmodern}
\usepackage{xcolor}
\usepackage{graphicx}
\usepackage{float}
\usepackage{booktabs}
\usepackage{longtable}
\usepackage{tabularx}
\usepackage{multirow}
\usepackage{array}
\usepackage{listings}
\usepackage{caption}
\usepackage{subcaption}
\usepackage{hyperref}
\usepackage{fancyhdr}
\usepackage{titlesec}
\usepackage{tcolorbox}
\usepackage{enumitem}
\usepackage{amsmath}
\usepackage{tikz}
\usepackage{pgfplots}
\pgfplotsset{compat=1.18}
\tcbuselibrary{skins,breakable}

%% ── Colors ────────────────────────────────────────────────────────────────────
\definecolor{primary}{HTML}{1E3A5F}
\definecolor{accent}{HTML}{2563EB}
\definecolor{lightblue}{HTML}{EFF6FF}
\definecolor{success}{HTML}{059669}
\definecolor{warning}{HTML}{D97706}
\definecolor{codeblue}{HTML}{0369A1}
\definecolor{codegray}{HTML}{F1F5F9}
\definecolor{codered}{HTML}{BE123C}
\definecolor{dockerblue}{HTML}{0DB7ED}
\definecolor{kafkaorange}{HTML}{E87722}
\definecolor{k8sblue}{HTML}{326CE5}
\definecolor{vaultgold}{HTML}{FFD700}

%% ── Hyperref Setup ────────────────────────────────────────────────────────────
\hypersetup{
  colorlinks=true, linkcolor=accent, urlcolor=accent, citecolor=accent,
  pdftitle={FinSight -- Intelligent Personal Finance Analytics Platform},
  pdfauthor={Aayank Singh},
  pdfsubject={SPE Major Project Report}
}

%% ── Header / Footer ───────────────────────────────────────────────────────────
\pagestyle{fancy}
\fancyhf{}
\fancyhead[L]{\small\textcolor{primary}{\textbf{FinSight} -- Intelligent Personal Finance Analytics Platform}}
\fancyhead[R]{\small\textcolor{primary}{SPE Major Project}}
\fancyfoot[C]{\thepage}
\renewcommand{\headrulewidth}{0.4pt}

%% ── Section Formatting ────────────────────────────────────────────────────────
\titleformat{\chapter}[display]
  {\normalfont\huge\bfseries\color{primary}}
  {\chaptertitlename\ \thechapter}{20pt}{\Huge}
\titleformat{\section}{\normalfont\Large\bfseries\color{primary}}{\thesection}{1em}{}
\titleformat{\subsection}{\normalfont\large\bfseries\color{accent}}{\thesubsection}{1em}{}
\titleformat{\subsubsection}{\normalfont\normalsize\bfseries\color{primary}}{\thesubsubsection}{1em}{}

%% ── Code Listings ─────────────────────────────────────────────────────────────
\lstset{
  basicstyle=\ttfamily\footnotesize,
  backgroundcolor=\color{codegray},
  frame=single, framerule=0pt, rulecolor=\color{codegray},
  breaklines=true, breakatwhitespace=true,
  keywordstyle=\color{codeblue}\bfseries,
  commentstyle=\color{success}\itshape,
  stringstyle=\color{codered},
  numberstyle=\tiny\color{gray},
  numbers=left, stepnumber=1,
  tabsize=2, showspaces=false, showstringspaces=false,
  captionpos=b
}

%% ── Custom Boxes ──────────────────────────────────────────────────────────────
\newtcolorbox{infobox}[1][]{
  enhanced, colback=lightblue, colframe=accent,
  fonttitle=\bfseries\color{white}, title=#1,
  boxrule=0.5pt, arc=4pt, left=6pt, right=6pt, top=6pt, bottom=6pt
}
\newtcolorbox{techbox}[1][]{
  enhanced, colback=codegray, colframe=primary,
  fonttitle=\bfseries\color{white}, title=#1, coltitle=white,
  boxrule=0.5pt, arc=4pt
}
\newtcolorbox{warningbox}[1][]{
  enhanced, colback=yellow!10, colframe=warning,
  fonttitle=\bfseries, title=#1,
  boxrule=0.5pt, arc=4pt
}

%% ── Document Start ────────────────────────────────────────────────────────────
\begin{document}

%% ── Title Page ────────────────────────────────────────────────────────────────
\begin{titlepage}
  \centering
  \vspace*{1.5cm}

  {\color{primary}\rule{\linewidth}{1.5pt}}\\[0.4cm]
  {\Huge\bfseries\color{primary} FinSight}\\[0.3cm]
  {\Large\color{accent} An Intelligent Personal Finance Analytics Platform}\\[0.2cm]
  {\large\color{primary} with MLOps-Driven Insights and Cloud-Native DevOps Automation}\\[0.1cm]
  {\color{primary}\rule{\linewidth}{1.5pt}}\\[1cm]

  \begin{tcolorbox}[enhanced, colback=lightblue, colframe=accent, width=0.85\linewidth, arc=6pt]
    \centering
    \textbf{Project Type:} Software and Platform Engineering (SPE) -- Major Project\\[4pt]
    \textbf{Domain:} Personal Finance $\times$ Machine Learning $\times$ Cloud-Native DevOps
  \end{tcolorbox}

  \vspace{1.5cm}

  % Screenshot placeholder
  \begin{tcolorbox}[enhanced, colback=gray!10, colframe=gray!50, width=0.9\linewidth,
      height=6cm, arc=6pt, valign=center, halign=center]
    \textcolor{gray!70}{\large\textit{[SCREENSHOT PLACEHOLDER: FinSight Dashboard Overview]}}
  \end{tcolorbox}

  \vfill

  \begin{tabular}{rl}
    \textbf{Student:}   & Aayank Singh Bhai \\[4pt]
    \textbf{Roll No.:}  & \textit{[Roll Number]} \\[4pt]
    \textbf{Course:}    & B.Tech. -- Computer Science \& Engineering \\[4pt]
    \textbf{Institute:} & \textit{[Institute Name]} \\[4pt]
    \textbf{Year:}      & 2025--2026 \\
  \end{tabular}

  \vspace{1cm}
  {\color{primary}\rule{\linewidth}{0.8pt}}
\end{titlepage}

%% ── Table of Contents ─────────────────────────────────────────────────────────
\tableofcontents
\newpage

%% ════════════════════════════════════════════════════════════════════════════
\chapter{Executive Summary}
%% ════════════════════════════════════════════════════════════════════════════

\begin{infobox}[Project at a Glance]
FinSight is a full-stack, production-grade, cloud-native personal finance platform built on a \textbf{polyglot microservices architecture}. It combines Java Spring Boot backend services, Python FastAPI ML services, a React 18 + TypeScript frontend, Apache Kafka for event streaming, HashiCorp Vault for secrets management, and a complete DevOps pipeline powered by Jenkins, Docker, Ansible, and Kubernetes -- all observed through the ELK Stack and Prometheus/Grafana.
\end{infobox}

FinSight enables users to upload bank statements, automatically parse and categorize transactions using machine learning, detect financial anomalies, forecast future spending, compute a financial stress score, and converse with an AI assistant grounded in their own financial data. The platform operationalizes MLOps best practices including model versioning (MLflow), drift detection, and automated retraining.

Key engineering highlights include:
\begin{itemize}[leftmargin=1.5em]
  \item \textbf{18 containerized microservices} (11 Java Spring Boot + 6 Python FastAPI + 1 React/Nginx frontend)
  \item \textbf{Kubernetes HPA} scaling ML inference pods from 2 to 8 replicas at 70\% CPU utilization
  \item \textbf{HashiCorp Vault} for zero-plaintext-secret policy across all services
  \item \textbf{ELK Stack} (Elasticsearch 8.11, Logstash, Kibana) for centralized structured log analytics
  \item \textbf{Kafka 3.x} event bus decoupling ingestion, categorization, anomaly detection, and notifications
  \item \textbf{Jenkins CI/CD} with 7-stage pipeline from checkout to Ansible-driven deployment
  \item \textbf{RS256 JWT} authentication with Google OAuth 2.0 support
  \item \textbf{Flyway} database migrations for zero-downtime schema evolution
\end{itemize}

%% ════════════════════════════════════════════════════════════════════════════
\chapter{Problem Statement}
%% ════════════════════════════════════════════════════════════════════════════

\section{Domain Context}

Personal finance management in India remains a largely manual and fragmented activity. Most individuals track expenses in spreadsheets, rely on bank statements received months after transactions, and have no proactive mechanism to understand their financial health. The few digital tools available either:

\begin{itemize}[leftmargin=1.5em]
  \item Apply rigid, rule-based categorization that fails on diverse bank statement formats,
  \item Offer static monthly summaries without predictive or anomaly-detection capability,
  \item Treat all users with a uniform behavioral model, ignoring individual spending DNA, or
  \item Expose user financial data to third-party advertisement networks.
\end{itemize}

\section{Core Problem}

\begin{tcolorbox}[enhanced, colback=primary!5, colframe=primary, arc=4pt, left=8pt, top=8pt]
\textbf{Problem:} How do we build a \emph{personalized, proactive, and privacy-preserving} financial intelligence platform that ingests raw multi-format bank statements, applies per-user machine learning models for categorization and anomaly detection, delivers real-time budget alerts, and remains horizontally scalable under production workloads -- all while maintaining a zero-trust security posture and full observability?
\end{tcolorbox}

\section{Specific Challenges Addressed}

\begin{longtable}{p{3.5cm} p{10.5cm}}
\toprule
\textbf{Challenge} & \textbf{FinSight's Approach} \\
\midrule
\endfirsthead
\multicolumn{2}{c}{\tablename\ \thetable{} -- continued}\\
\toprule
\textbf{Challenge} & \textbf{FinSight's Approach} \\
\midrule
\endhead
\bottomrule
\endlastfoot

Bank format heterogeneity & Intelligent column detection in \texttt{ingestion-service} using heuristic matching across CSV/XLS/PDF formats \\[4pt]
Per-user behavioral modeling & LSTM autoencoder in \texttt{anomaly-detection-service} learns individual spending DNA, not a population average \\[4pt]
Secrets sprawl & HashiCorp Vault with Kubernetes Secrets bridge -- zero plaintext credentials in any config file \\[4pt]
ML inference scalability & Kubernetes HPA on \texttt{anomaly-detection-service} and \texttt{categorization-service} (2--8 replicas, CPU \textgreater 70\%) \\[4pt]
Real-time budget alerting & Kafka \texttt{budget.alerts} topic $\rightarrow$ \texttt{notification-service} $\rightarrow$ WebSocket/STOMP push to frontend \\[4pt]
Log observability at scale & ELK Stack: structured JSON logs from all 18 services $\rightarrow$ Logstash $\rightarrow$ Elasticsearch $\rightarrow$ Kibana dashboards \\[4pt]
Zero-downtime deployments & Kubernetes RollingUpdate (\texttt{maxUnavailable: 0, maxSurge: 1}) across all service deployments \\[4pt]
Model drift & KL divergence computed in \texttt{model-training-service}; admin-triggerable retraining via Kafka \\[4pt]
\end{longtable}

%% ════════════════════════════════════════════════════════════════════════════
\chapter{Technology Stack}
%% ════════════════════════════════════════════════════════════════════════════

\section{Stack Overview}

\begin{figure}[H]
  \centering
  \begin{tcolorbox}[enhanced, colback=gray!10, colframe=gray!50,
      width=\linewidth, height=5cm, arc=6pt, valign=center, halign=center]
    \textcolor{gray!70}{\large\textit{[SCREENSHOT PLACEHOLDER: Technology Stack Diagram]}}
  \end{tcolorbox}
  \caption{FinSight Full Technology Stack}
\end{figure}

\section{Frontend}

\begin{techbox}[Frontend -- React 18 + TypeScript + Vite]
\begin{itemize}[leftmargin=1em, noitemsep]
  \item \textbf{React 18} with TypeScript -- component framework with strict typing
  \item \textbf{Vite} -- HMR-enabled build tool replacing Create React App
  \item \textbf{Tailwind CSS} -- utility-first responsive styling
  \item \textbf{Recharts} -- spending charts, trend lines, category breakdowns
  \item \textbf{TanStack Query (React Query)} -- server-state management and API caching
  \item \textbf{Zustand} -- lightweight client-state (auth session, user profile)
  \item \textbf{React Router v6} -- client-side routing with protected routes
  \item \textbf{Axios} -- HTTP client with JWT interceptors
  \item \textbf{@react-oauth/google} -- Google OAuth 2.0 frontend flow
\end{itemize}
\end{techbox}

\section{Backend -- Java Spring Boot Microservices}

All Java services use \textbf{Java 21 + Spring Boot 3.x}, each with its own \texttt{Dockerfile}, database connection, Flyway migration history, and Spring Actuator health endpoints.

\begin{longtable}{p{3.8cm} p{5cm} p{5.2cm}}
\toprule
\textbf{Service} & \textbf{Responsibility} & \textbf{Key Dependencies} \\
\midrule
\endfirsthead
\toprule
\textbf{Service} & \textbf{Responsibility} & \textbf{Key Dependencies} \\
\midrule
\endhead
\bottomrule
\endlastfoot
\texttt{api-gateway-service} & Auth enforcement, routing, rate limiting & Spring Cloud Gateway, Spring Security, jjwt \\[3pt]
\texttt{user-service} & Registration, login, OAuth, profile & Spring Security, OAuth2 Client, Spring Data JPA \\[3pt]
\texttt{ingestion-service} & File upload, parsing, Kafka publish & Spring Kafka, Apache POI, PDFBox, OpenCSV \\[3pt]
\texttt{transaction-service} & CRUD, rules engine, deduplication & Spring Data JPA, Spring Kafka consumer \\[3pt]
\texttt{budget-service} & Budget limits, real-time tracking, threshold alerts & Spring Data JPA, Spring Kafka \\[3pt]
\texttt{notification-service} & In-app alerts, inbox, WebSocket push & Spring Kafka consumer, WebSocket (STOMP) \\[3pt]
\texttt{admin-service} & User management, platform metrics, ML config & Spring Data JPA, Spring Security \\[3pt]
\texttt{chat-service} & Query intake, ML orchestration, Gemini AI & Spring WebFlux, WebClient, Gemini API \\[3pt]
\end{longtable}

\section{Backend -- Python FastAPI ML Microservices}

All ML services use \textbf{Python 3.11 + FastAPI}. Separation from Spring Boot is deliberate -- Python owns the data science ecosystem; Spring Boot owns user-facing transactional logic.

\begin{longtable}{p{4.2cm} p{4.8cm} p{5cm}}
\toprule
\textbf{Service} & \textbf{Responsibility} & \textbf{Key Libraries} \\
\midrule
\endfirsthead
\toprule
\textbf{Service} & \textbf{Responsibility} & \textbf{Key Libraries} \\
\midrule
\endhead
\bottomrule
\endlastfoot
\texttt{categorization-service} & ML transaction categorization & scikit-learn, XGBoost, FastAPI, kafka-python \\[3pt]
\texttt{anomaly-detection-service} & Per-user LSTM anomaly scorer & PyTorch, numpy, pandas, scikit-learn, kafka-python \\[3pt]
\texttt{forecasting-service} & Category-wise spend forecasting & pandas, SQLAlchemy, psycopg2, PyJWT, kafka-python \\[3pt]
\texttt{parsing-ml-service} & Intelligent bank statement parsing & pandas, pdfplumber, openpyxl, spaCy \\[3pt]
\texttt{stress-score-service} & Composite financial stress scorer & scikit-learn, pandas, FastAPI \\[3pt]
\texttt{model-training-service} & Retraining, drift detection, MLflow & MLflow, PyTorch, scikit-learn, scipy \\[3pt]
\endlongtable}

\section{Data Stores}

\begin{longtable}{p{2.5cm} p{4.5cm} p{7cm}}
\toprule
\textbf{Store} & \textbf{Used By} & \textbf{Rationale} \\
\midrule
\endfirsthead
\toprule
\textbf{Store} & \textbf{Used By} & \textbf{Rationale} \\
\midrule
\endhead
\bottomrule
\endlastfoot
PostgreSQL 15 & user-service, transaction-service, budget-service, notification-service, chat-service, admin-service & Relational, ACID-compliant structured financial data; schema-per-service isolation; Flyway migrations \\[3pt]
MongoDB 7.0 & ingestion-service & Flexible document model for raw file metadata and parsing job state \\[3pt]
Redis 7 & api-gateway (rate limiting, token cache), chat-service (session), simulation-service (cache) & Sub-millisecond reads, ephemeral data, token revocation set \\[3pt]
\endlongtable}

\section{Messaging -- Apache Kafka 3.x}

\begin{techbox}[Kafka Topics]
\begin{tabular}{ll}
\texttt{transactions.ingested}      & Ingestion $\rightarrow$ Transaction service \& Anomaly detection \\
\texttt{transactions.categorized}   & Categorization service $\rightarrow$ Transaction service \\
\texttt{anomaly.alerts}             & Anomaly detection $\rightarrow$ Notification service \\
\texttt{budget.alerts}              & Budget service $\rightarrow$ Notification service \\
\texttt{model.retrain.trigger}      & Admin service $\rightarrow$ Model training service \\
\end{tabular}
\end{techbox}

\section{DevOps \& Infrastructure Toolchain}

\begin{longtable}{p{3.5cm} p{10.5cm}}
\toprule
\textbf{Tool} & \textbf{Role in FinSight} \\
\midrule
\endfirsthead
\toprule
\textbf{Tool} & \textbf{Role in FinSight} \\
\midrule
\endhead
\bottomrule
\endlastfoot
Git + GitHub & Version control, webhook trigger for Jenkins \\
Jenkins & 7-stage CI/CD pipeline (Checkout $\rightarrow$ Build $\rightarrow$ Test $\rightarrow$ Docker $\rightarrow$ Push $\rightarrow$ Ansible $\rightarrow$ K8s) \\
Docker + Compose & Containerization of all 18 services; local dev environment \\
Docker Hub & Container image registry (\texttt{aayanksinghai/finsight-*}) \\
Kubernetes (Minikube) & Production-like orchestration in \texttt{finsight} namespace \\
Kubernetes HPA & Auto-scaling ML inference pods 2--8 replicas at CPU $>$ 70\% \\
Ansible & Configuration management; deploy playbook \\
HashiCorp Vault & Zero-plaintext secrets injection at runtime \\
Prometheus + Grafana & Metrics scraping and infrastructure dashboards \\
ELK Stack & Centralized log aggregation and Kibana analytics \\
MLflow & ML experiment tracking, model versioning, model registry \\
\endlongtable}

%% ════════════════════════════════════════════════════════════════════════════
\chapter{System Architecture}
%% ════════════════════════════════════════════════════════════════════════════

\section{High-Level Architecture}

\begin{figure}[H]
  \centering
  \begin{tcolorbox}[enhanced, colback=gray!10, colframe=gray!50,
      width=\linewidth, height=9cm, arc=6pt, valign=center, halign=center]
    \textcolor{gray!70}{\Large\textit{[SCREENSHOT PLACEHOLDER: Full System Architecture Diagram]}}\\[6pt]
    \textcolor{gray!60}{\small React Frontend $\rightarrow$ API Gateway $\rightarrow$ [Spring Boot Services] $\rightarrow$ [Python ML Services]}\\
    \textcolor{gray!60}{\small Kafka event bus $\leftrightarrow$ Notification service $\rightarrow$ WebSocket $\rightarrow$ Frontend}
  \end{tcolorbox}
  \caption{FinSight System Architecture -- Microservices, Messaging, and Observability}
\end{figure}

\section{Architectural Layers}

\subsection{Layer 1 -- Presentation (React SPA)}
The frontend is a single-page application served by Nginx inside a Docker container. All API requests are proxied through the API Gateway on port \texttt{8090}. The React app uses Zustand for auth state and TanStack Query for server-state caching. WebSocket connections for real-time notifications are established via SockJS/STOMP directly to the \texttt{notification-service}.

\subsection{Layer 2 -- API Gateway (Spring Cloud Gateway)}
The API Gateway (\texttt{port 8090}, \texttt{NodePort 30090} in K8s) is the single entry point for all HTTP traffic. It enforces:
\begin{itemize}[leftmargin=1.5em, noitemsep]
  \item \textbf{JWT validation} via RS256 public key -- all protected routes validated before forwarding
  \item \textbf{Rate limiting} -- 100 requests per minute per authenticated user (Redis token bucket)
  \item \textbf{CORS policy} for frontend origin
  \item \textbf{Route configuration} to downstream services by path prefix
\end{itemize}

\subsection{Layer 3 -- Java Spring Boot Services}
Eight stateless Spring Boot services handle all user-facing transactional logic. Each service:
\begin{itemize}[leftmargin=1.5em, noitemsep]
  \item Exposes \texttt{/actuator/health/liveness} and \texttt{/actuator/health/readiness} for K8s probes
  \item Publishes structured JSON logs to Logstash via the ELK pipeline
  \item Uses Flyway for database migration with per-service migration history tables
  \item Injects secrets from K8s Secrets (populated from Vault) at runtime
\end{itemize}

\subsection{Layer 4 -- Python FastAPI ML Services}
Six FastAPI services own the data-science layer. They communicate with Spring Boot via:
\begin{itemize}[leftmargin=1.5em, noitemsep]
  \item \textbf{Synchronous REST}: Spring Boot WebClient calls Python ML endpoints for inference
  \item \textbf{Asynchronous Kafka}: categorization-service and anomaly-detection-service consume \texttt{transactions.ingested} and publish results back
\end{itemize}

\subsection{Layer 5 -- Event Bus (Apache Kafka)}
Kafka is the asynchronous backbone decoupling ingestion from downstream processing. The \texttt{confluentinc/cp-kafka:7.5.0} image is used with Zookeeper. Key design: ingestion events are published by the ingestion-service; Python ML services consume and enrich; notification-service delivers real-time alerts to users.

\subsection{Layer 6 -- Observability (ELK + Prometheus/Grafana)}
Every microservice emits structured JSON logs to Logstash on port 5044. Logstash forwards to Elasticsearch 8.11. Kibana provides pre-built dashboards for ingestion rate, ML inference latency, error rate per service, and anomaly alert volume. Spring Boot Actuator and Micrometer expose Prometheus-format metrics scraped by Prometheus; Grafana visualizes CPU, memory, and request throughput per pod.

\section{Inter-Service Communication Patterns}

\begin{longtable}{p{3cm} p{5cm} p{6cm}}
\toprule
\textbf{Pattern} & \textbf{Used Between} & \textbf{Implementation} \\
\midrule
\endfirsthead
\toprule
\textbf{Pattern} & \textbf{Used Between} & \textbf{Implementation} \\
\midrule
\endhead
\bottomrule
\endlastfoot
Synchronous REST & React $\rightarrow$ Gateway $\rightarrow$ Spring Boot $\rightarrow$ Python ML & Spring WebFlux WebClient (reactive, non-blocking) \\[3pt]
Asynchronous Events & Spring Boot $\rightarrow$ Kafka $\rightarrow$ Spring Boot / Python & spring-kafka / kafka-python consumers \\[3pt]
WebSocket (STOMP) & notification-service $\rightarrow$ React frontend & SockJS over WebSocket, JWT auth on handshake \\[3pt]
\endlongtable}

\section{Authentication Flow}

\begin{enumerate}[leftmargin=1.5em]
  \item \textbf{Standard login:} User submits credentials $\rightarrow$ \texttt{user-service} validates bcrypt hash $\rightarrow$ issues RS256-signed JWT $\rightarrow$ stored in \texttt{httpOnly} cookie on frontend.
  \item \textbf{Google OAuth:} React uses \texttt{@react-oauth/google} $\rightarrow$ sends \texttt{id\_token} to \texttt{user-service} $\rightarrow$ validated via Google's JWKS endpoint $\rightarrow$ JWT issued (new account created or existing linked on email match).
  \item \textbf{Subsequent requests:} JWT validated at \texttt{api-gateway-service} via Spring Security filter chain before any route is forwarded.
  \item \textbf{Token refresh:} Refresh token sessions persisted in \texttt{refresh\_token\_sessions} table; access tokens tracked in \texttt{revoked\_access\_tokens} for invalidation on logout.
\end{enumerate}

%% ════════════════════════════════════════════════════════════════════════════
\chapter{Microservices Specification \& DB Schemas}
%% ════════════════════════════════════════════════════════════════════════════

\section{Service Breakdown}

The FinSight platform is decomposed into 14 distinct core backend microservices. The separation between Java Spring Boot and Python FastAPI is strictly enforced based on the nature of workload (user-facing transactional vs. intensive machine learning inference).

\subsection{Java Spring Boot Microservices}

\subsubsection{1. API Gateway Service (\texttt{api-gateway-service})}
Acts as the single front door for all client traffic. Built on Spring Cloud Gateway, it implements global filters for authorization. It intercepts incoming HTTP requests, extracts the RS256 JWT, validates its signature against the public key, and enforces rate limiting (100 requests/min/user) using a Redis token bucket algorithm.

\subsubsection{2. User Service (\texttt{user-service})}
Manages user onboarding, authentication, profile metadata, and session lifecycle. Supports email/password registration with bcrypt hashing (cost factor 12) and Google OAuth 2.0.
\begin{itemize}[leftmargin=1.5em, noitemsep]
  \item \textbf{DB Schema Table (\texttt{user\_credentials}):} \texttt{id} (UUID), \texttt{email} (VARCHAR 320, UNIQUE), \texttt{password\_hash} (VARCHAR 255), \texttt{full\_name}, \texttt{city}, \texttt{age\_group}, \texttt{monthly\_income}, \texttt{role} (VARCHAR 20), \texttt{deactivated\_at}.
  \item \textbf{Token Tracking Tables:} \texttt{refresh\_token\_sessions} (\texttt{id}, \texttt{token\_jti}, \texttt{user\_email}, \texttt{expires\_at}) and \texttt{revoked\_access\_tokens}.
\end{itemize}

\subsubsection{3. Ingestion Service (\texttt{ingestion-service})}
Orchestrates the ingestion of bank statement files (CSV, XLS, PDF). It stores uploaded file metadata and asynchronous parse job documents in MongoDB. Once parsed into canonical rows, it publishes \texttt{transactions.ingested} events to Kafka.

\subsubsection{4. Transaction Service (\texttt{transaction-service})}
The core ledger of the platform. Consumes ingestion events, performs deduplication using a cryptographic hash (\texttt{content\_hash} generated from date + amount + description), and manages categories.
\begin{itemize}[leftmargin=1.5em, noitemsep]
  \item \textbf{DB Schema Table (\texttt{transactions}):} \texttt{id} (UUID), \texttt{content\_hash} (VARCHAR 64, UNIQUE), \texttt{owner\_email}, \texttt{account\_id}, \texttt{job\_id}, \texttt{source\_file\_name}, \texttt{source\_bank}, \texttt{occurred\_at}, \texttt{raw\_description}, \texttt{normalized\_merchant}, \texttt{category\_id}, \texttt{amount} (NUMERIC 15,2), \texttt{type} (DEBIT/CREDIT), \texttt{currency}, \texttt{tags} (TEXT[]), \texttt{is\_anomaly}.
\end{itemize}

\subsubsection{5. Budget Service (\texttt{budget-service})}
Enables users to set monthly spending caps per category. Consumes transaction events from Kafka to update running totals in real time. If spend breaches 80\% or 100\% of a budget limit, it publishes events to \texttt{budget.alerts}.
\begin{itemize}[leftmargin=1.5em, noitemsep]
  \item \textbf{DB Schema Table (\texttt{budgets}):} \texttt{id} (UUID), \texttt{owner\_email}, \texttt{category\_id}, \texttt{category\_name}, \texttt{limit\_amount}, \texttt{current\_spend}, \texttt{month\_year} (YYYY-MM).
\end{itemize}

\subsubsection{6. Notification Service (\texttt{notification-service})}
Processes \texttt{budget.alerts} and \texttt{anomaly.alerts} from Kafka. It persists alerts in a notification inbox and pushes real-time updates over WebSocket/STOMP to active client sessions.
\begin{itemize}[leftmargin=1.5em, noitemsep]
  \item \textbf{DB Schema Table (\texttt{notifications}):} \texttt{id} (UUID), \texttt{owner\_email}, \texttt{title}, \texttt{type}, \texttt{message}, \texttt{is\_read}, \texttt{created\_at}.
  \item \textbf{DB Schema Table (\texttt{notification\_preferences}):} \texttt{owner\_email}, \texttt{type}, \texttt{enabled}.
\end{itemize}

\subsubsection{7. Chat Service (\texttt{chat-service})}
Provides a conversational AI interface using the Gemini API. It intercepts natural language questions, queries internal microservices (transaction, budget, forecasting) via REST WebClient for grounding context, and synthesizes personalized responses.
\begin{itemize}[leftmargin=1.5em, noitemsep]
  \item \textbf{DB Schema Tables:} \texttt{chat\_sessions} (\texttt{id}, \texttt{owner\_email}), \texttt{chat\_messages} (\texttt{id}, \texttt{session\_id}, \texttt{role}, \texttt{content}, \texttt{intent}), \texttt{chat\_ratings} (\texttt{message\_id}, \texttt{rating}).
\end{itemize}

\subsubsection{8. Admin Service (\texttt{admin-service})}
Provides administrative oversight. Exposes endpoints for user deactivation, platform telemetry aggregation, and dynamic ML threshold adjustment.
\begin{itemize}[leftmargin=1.5em, noitemsep]
  \item \textbf{DB Schema Table (\texttt{admin\_config}):} \texttt{config\_key} (PRIMARY KEY), \texttt{config\_value}, \texttt{description}. Stores parameters like \texttt{anomaly.sensitivity} and \texttt{stress.weight.*}.
\end{itemize}

\subsection{Python FastAPI ML Microservices}

\subsubsection{9. Parsing ML Service (\texttt{parsing-ml-service})}
Utilizes NLP and heuristic column mapping to intelligently extract transaction records from unstructured or multi-column bank statements without requiring hardcoded templates.

\subsubsection{10. Categorization Service (\texttt{categorization-service})}
Employs an XGBoost classifier trained on merchant names and transaction descriptions to classify transactions into standard financial categories. Listens to Kafka for unclassified transactions and publishes categorized outcomes.

\subsubsection{11. Anomaly Detection Service (\texttt{anomaly-detection-service})}
Runs a PyTorch-based LSTM autoencoder. For each user, it builds an individual spending DNA baseline. It scores incoming transactions in near real-time (target latency $<$ 200ms) and flags anomalous behavior.

\subsubsection{12. Forecasting Service (\texttt{forecasting-service})}
Uses time-series forecasting algorithms (LSTM/ARIMA) over historical transaction data to predict next month's spending across categories, enabling proactive budget vs. prediction comparisons.

\subsubsection{13. Stress Score Service (\texttt{stress-score-service})}
Computes a monthly composite financial stress score (0--100) combining spend-to-income ratios, savings rates, recurring payment burden, and budget adherence.

\subsubsection{14. Simulation Service (\texttt{simulation-service})}
Powers the interactive "What-If" simulator in the frontend. Evaluates hypothetical spending modifications in real-time by chaining forecasting and stress scoring models without mutating production state.

\subsubsection{15. Model Training Service (\texttt{model-training-service})}
Manages scheduled and admin-triggered model retraining. Computes KL divergence to detect data drift and integrates with MLflow for model tracking and artifact versioning.

%% ════════════════════════════════════════════════════════════════════════════
\chapter{Cloud-Native Infrastructure \& DevOps Automation}
%% ════════════════════════════════════════════════════════════════════════════

\section{Docker Containerization}

Every microservice in FinSight is fully containerized. A multi-stage build approach is adopted for Java Spring Boot services using Maven to compile and package the application, resulting in lean runtime images based on Eclipse Temurin 21 Alpine. Python services use \texttt{uvicorn} with lean Python 3.11 slim base images.

\begin{figure}[H]
  \centering
  \begin{tcolorbox}[enhanced, colback=gray!10, colframe=gray!50, width=0.85\linewidth,
      height=5cm, arc=6pt, valign=center, halign=center]
    \textcolor{gray!70}{\large\textit{[SCREENSHOT PLACEHOLDER: Docker Compose Infrastructure Setup]}}
  \end{tcolorbox}
  \caption{Docker Compose Multi-Container Deployment}
\end{figure}

\section{HashiCorp Vault Integration}

To achieve a zero-trust security posture, no plaintext passwords or API keys exist in source code or configuration files. HashiCorp Vault is deployed as the central secrets engine.

\begin{itemize}[leftmargin=1.5em]
  \item \textbf{Secret Engine Path:} Secrets are stored under \texttt{secret/data/finsight}.
  \item \textbf{Kubernetes Injection:} In Kubernetes, a bridge mechanism (External Secrets Operator / init-containers) populates Kubernetes Secrets (\texttt{finsight-secrets}) directly from Vault.
  \item \textbf{Runtime Binding:} Deployments inject these secrets as environment variables into pods (e.g., \texttt{DB\_PASSWORD}, \texttt{JWT\_PRIVATE\_KEY}, \texttt{GEMINI\_API\_KEY}).
\end{itemize}

\section{Kubernetes Architecture \& Configuration}

FinSight is deployed into a dedicated Kubernetes namespace (\texttt{finsight}). The architecture leverages Deployments, StatefulSets, Services, ConfigMaps, and Secrets.

\begin{figure}[H]
  \centering
  \begin{tcolorbox}[enhanced, colback=gray!10, colframe=gray!50, width=0.9\linewidth,
      height=6cm, arc=6pt, valign=center, halign=center]
    \textcolor{gray!70}{\large\textit{[SCREENSHOT PLACEHOLDER: Kubernetes Cluster Workloads \& Pods]}}
  \end{tcolorbox}
  \caption{Kubernetes Pods, Services, and StatefulSets in \texttt{finsight} Namespace}
\end{figure}

\subsection{Stateful Data Storage (PostgreSQL)}
To ensure reliable persistence and predictable DNS resolution, PostgreSQL is deployed as a \texttt{StatefulSet} with a Headless Service (\texttt{clusterIP: None}). It attaches a \texttt{PersistentVolumeClaim} requesting 5Gi of ReadWriteOnce storage.

\subsection{Probes \& Rolling Updates}
All Deployments configure strict liveness and readiness probes pointing to \texttt{/health} or Actuator endpoints. To satisfy the zero-downtime requirement, the Deployment strategy is explicitly configured:
\begin{lstlisting}[language=YAML, caption=Kubernetes RollingUpdate Strategy]
strategy:
  type: RollingUpdate
  rollingUpdate:
    maxUnavailable: 0
    maxSurge: 1
\end{lstlisting}

\section{Horizontal Pod Autoscaling (HPA)}

Machine learning inference (especially LSTM scoring) is highly CPU-intensive. To maintain SLA under traffic spikes without over-provisioning infrastructure, \texttt{HorizontalPodAutoscaler} resources are bound to \texttt{anomaly-detection-service} and \texttt{categorization-service}.

\begin{lstlisting}[language=YAML, caption=HPA Configuration for ML Inference Services]
apiVersion: autoscaling/v2
kind: HorizontalPodAutoscaler
metadata:
  name: anomaly-detection-service-hpa
  namespace: finsight
spec:
  scaleTargetRef:
    apiVersion: apps/v1
    kind: Deployment
    name: anomaly-detection-service
  minReplicas: 2
  maxReplicas: 8
  metrics:
    - type: Resource
      resource:
        name: cpu
        target:
          type: Utilization
          averageUtilization: 70
\end{lstlisting}
When average CPU utilization across pods exceeds 70\%, Kubernetes automatically scales the deployment up to 8 replicas.

\section{Ansible Automation}

Infrastructure provisioning and application deployment across environments are automated using Ansible. The setup mandates four distinct roles:
\begin{enumerate}[leftmargin=1.5em]
  \item \textbf{\texttt{common}:} Prepares base OS, installs prerequisite libraries, and configures Docker.
  \item \textbf{\texttt{kafka}:} Deploys and configures Apache Kafka and Zookeeper clusters.
  \item \textbf{\texttt{kubernetes}:} Applies Kubernetes namespace, ConfigMaps, Secrets, and manifest files idempotently.
  \item \textbf{\texttt{monitoring}:} Deploys Prometheus, Grafana, and ELK Stack agents.
\end{enumerate}

\section{Jenkins CI/CD Pipeline}

Continuous integration and deployment are driven by a robust \texttt{Jenkinsfile} triggered automatically via GitHub webhooks on push to the \texttt{main} branch.

\begin{figure}[H]
  \centering
  \begin{tcolorbox}[enhanced, colback=gray!10, colframe=gray!50, width=0.9\linewidth,
      height=5cm, arc=6pt, valign=center, halign=center]
    \textcolor{gray!70}{\large\textit{[SCREENSHOT PLACEHOLDER: Jenkins CI/CD Pipeline Execution]}}
  \end{tcolorbox}
  \caption{Jenkins 7-Stage Automated CI/CD Pipeline}
\end{figure}

The pipeline executes the following sequential stages:
\begin{enumerate}[leftmargin=1.5em, noitemsep]
  \item \textbf{Checkout:} Pulls the latest source code from GitHub.
  \item \textbf{Build Backend:} Patches internal service discovery URLs and compiles Java microservices using Maven (\texttt{mvn package}).
  \item \textbf{Build Frontend:} Compiles the React/TypeScript frontend via Vite (\texttt{npm run build}).
  \item \textbf{Test:} Executes unit and integration test suites, generating JUnit/Surefire XML reports.
  \item \textbf{Build \& Push Docker Images:} Builds Docker images for all 12 services and pushes them to Docker Hub tagged as \texttt{latest}.
  \item \textbf{Deploy with Ansible:} Executes the Ansible deploy playbook (\texttt{deploy.yml}) to pull images and update container orchestrators.
  \item \textbf{Post-Build Cleanup:} Prunes dangling images to free worker node disk space.
\end{enumerate}

%% ════════════════════════════════════════════════════════════════════════════
\chapter{Observability \& Logging Pipeline}
%% ════════════════════════════════════════════════════════════════════════════

\section{ELK Stack Architecture}

A comprehensive centralized logging pipeline is established using Elasticsearch, Logstash, and Kibana. 

\begin{figure}[H]
  \centering
  \begin{tcolorbox}[enhanced, colback=gray!10, colframe=gray!50, width=0.9\linewidth,
      height=6cm, arc=6pt, valign=center, halign=center]
    \textcolor{gray!70}{\large\textit{[SCREENSHOT PLACEHOLDER: Kibana Log Analytics Dashboard]}}
  \end{tcolorbox}
  \caption{Kibana Centralized Telemetry and Log Analytics Dashboard}
\end{figure}

\begin{itemize}[leftmargin=1.5em]
  \item \textbf{Log Emission:} All Spring Boot and FastAPI microservices emit structured JSON logs to TCP port 5044.
  \item \textbf{Logstash Aggregation:} Logstash ingests JSON streams, applies grok/filtering rules, and indexes records into Elasticsearch.
  \item \textbf{Kibana Visualizations:} Pre-built Kibana dashboards monitor real-time transaction ingestion rates, ML inference latencies, anomaly trigger volumes, and per-service error rates.
  \item \textbf{Distributed Tracing:} A \texttt{correlation\_id} is injected into request headers at the API Gateway and propagated across HTTP/Kafka hops for end-to-end trace reconstruction.
\end{itemize}

\section{Prometheus \& Grafana Telemetry}

System and JVM telemetry are scraped by Prometheus and visualized in Grafana.
\begin{itemize}[leftmargin=1.5em]
  \item \textbf{Metrics Exposition:} Spring Boot microservices utilize Micrometer to expose Prometheus metrics at \texttt{/actuator/prometheus}.
  \item \textbf{Grafana Dashboards:} Real-time monitoring panels visualize CPU throttling, memory consumption, JVM heap usage, and HTTP request throughput per pod.
\end{itemize}

%% ════════════════════════════════════════════════════════════════════════════
\chapter{Innovative Solutions Implemented}
%% ════════════════════════════════════════════════════════════════════════════

\section{Personalized Anomaly Detection (LSTM Autoencoder)}
Standard financial platforms employ static rule-based thresholds (e.g., flagging transactions over ₹50,000). FinSight utilizes a PyTorch LSTM autoencoder that learns the temporal sequence of an individual user's spending habits (merchant frequency, typical transaction times, amount distribution). A ₹500 transaction at 11 PM at a petrol pump is considered normal for a sales executive but flagged as highly anomalous for a student.

\section{Composite Financial Stress Score}
FinSight computes an advanced monthly financial health metric (0--100) aggregating five critical financial vectors:
\begin{equation}
\text{Stress Score} = w_1\left(\frac{\text{Spend}}{\text{Income}}\right) + w_2(\text{Savings Rate}) + w_3(\text{EMI Burden}) + w_4(\text{Discretionary Growth}) + w_5(\text{Budget Breach})
\end{equation}
Where weights $w_i$ are dynamically configurable via the Admin Service. Changes greater than 10 points month-on-month trigger proactive alerts.

\section{What-If Financial Simulator}
An interactive simulation engine allowing users to manipulate sliders (e.g., "Reduce dining spend by ₹2,000" or "Add new EMI of ₹5,000"). The frontend calls the \texttt{simulation-service}, which dynamically evaluates the impact on next month's cash flow and stress score without modifying live ledger state.

\begin{figure}[H]
  \centering
  \begin{tcolorbox}[enhanced, colback=gray!10, colframe=gray!50, width=0.85\linewidth,
      height=5cm, arc=6pt, valign=center, halign=center]
    \textcolor{gray!70}{\large\textit{[SCREENSHOT PLACEHOLDER: What-If Simulator UI]}}
  \end{tcolorbox}
  \caption{What-If Financial Simulator Interface}
\end{figure}

\section{Spending Mood Correlation}
Users optionally log daily mood check-ins (Happy, Stressed, Sad). Over time, machine learning correlates emotional states with transaction categories, generating behavioral insights such as: \textit{"You spend 60\% more on online food delivery on days rated as Stressed."}

\section{Contextual AI Chat Assistant}
Rather than providing static FAQ responses, the AI Chat Assistant utilizes Gemini LLM grounded in real user data. When a user queries \textit{"Why did I overspend in March?"}, the service orchestrates internal calls to the anomaly service and budget service, synthesizing a precise, data-backed natural language explanation.

%% ════════════════════════════════════════════════════════════════════════════
\chapter{Conclusion \& Roadmap}
%% ════════════════════════════════════════════════════════════════════════════

FinSight successfully demonstrates the convergence of modern personal finance analytics, machine learning, and enterprise DevOps automation. By decoupling services via Kafka, securing secrets with Vault, ensuring observability through ELK/Prometheus, and implementing zero-downtime Kubernetes deployments, the platform establishes a highly scalable and robust architecture. Future development roadmap includes multi-bank account aggregation via Account Aggregator (AA) frameworks and automated tax harvesting optimization.

\end{document}
